\documentclass[11pt]{article}
\usepackage[a4paper,margin=28mm]{geometry}
\usepackage[T1]{fontenc}
\usepackage{lmodern,amsmath,amssymb,amsthm,mathtools,hyperref}
\newcommand{\ddc}{dd^c}
\newcommand{\PSH}{\operatorname{PSH}}
\newcommand{\MA}{\operatorname{MA}}
\newcommand{\E}{\mathcal E}
\newcommand{\Pone}{\mathbf P^1}
\newtheorem{theorem}{Theorem}
\newtheorem{lemma}[theorem]{Lemma}
\theoremstyle{remark}\newtheorem{remark}[theorem]{Remark}
\title{Selection with moving logarithmic ordinary models\newline
on Hermitian products}
\author{}\date{}
\begin{document}\maketitle
\begin{abstract}
For a family of logarithmic divisor models on Hermitian products of any dimension,
the greatest positive-exponent solutions select a critical potential as the
exponent tends to zero. The normalized limit retains the actual ordinary
reference type. A contraction on the curve factor gives uniform control
while the reference coefficients, background and continuous density vary
at independent rates. The selected potential has an order characterization
in the separated sector. The scalar-centering error has an explicit limit,
provided centering uses each parameter's critical scalar. Centering by the
limiting scalar can instead introduce a divergent constant.
The critical coefficient is maximal among all subsolutions in the prescribed
one-sided ordinary singularity class.
\end{abstract}

Use $\ddc=i\partial\bar\partial/(2\pi)$ and ordinary wedge powers.
Potentials have their genuine canonical upper representatives. All products
are the non-pluripolar products with smooth-background corrections.
The assertions concern the specified singular family, not arbitrary model
references or uniqueness of arbitrary nonclosed critical potentials.

\section{The geometry and the selected potential}

Let $(Y,\kappa)$ be any connected compact Hermitian $k$-fold, $k\geq1$,
normalized by $\int_Y\kappa^k=1$. On $X=\Pone_x\times Y$ put $N=k+1$,
take the mass-one Fubini--Study form $\omega_1$ and use
$dW=\omega_1\wedge\kappa^k$. Pullbacks are implicit; $\kappa$ need not be closed.
Fix $a,d>0$, a smooth positive nonconstant $b=b(x)$, and continuous $H=H(x)>0$.
Set
\[
 \theta=(a+d)\omega_1+b\kappa,\quad
 \eta=a\omega_1+b\kappa,\quad f=N b^kH,\quad
 m=\int_{\Pone}H\omega_1,\quad c_0=\frac am.
\]
Both backgrounds are strictly positive and genuinely nonclosed. Indeed the
component of $\ddc\theta$ with horizontal real degree two is
$(\ddc b)\wedge\kappa\ne0$; the other components have different horizontal degrees.
With $D=\{Z_0=0\}\times Y$, write
\[
 \ell=\frac{|Z_0|^2}{|Z_0|^2+|Z_1|^2},\qquad
 \psi=d\log\ell,\qquad
 \theta+\ddc\psi=\eta+d[D].
\]
The canonical value of $\psi$ on $D$ is $-\infty$.

For $K=a+d$, $R=a/d$ and
$C_*=K\log(K/d)-a\log(a/d)$, define in $z=Z_1/Z_0$
\begin{equation}\label{eq:model}
 \phi_{a,d}=\begin{cases}
 0,& |z|^2\leq R,\\
 a\log|z|^2+C_*-K\log(1+|z|^2),& |z|^2>R,
 \end{cases}\qquad \phi_{a,d}|_D=-\infty.
\end{equation}
We verify its ordinary model property below. In particular
$\psi\leq\phi_{a,d}\leq\psi+C_*$.
Write $v\preceq\phi$ for $v\leq\phi+C$ and $v\sim\phi$ for mutual
bounded comparisons, including their common canonical poles.

Let
\[
 \begin{aligned}
 G(x,z)&=\log d_{\rm chord}(x,z)^2+1,
 &w_0&=G*\left(\frac{aH}{m}-a\right),\\
 S_0&=\sup_X(\psi+w_0),
 &d\nu_0&=\frac Hm\,dW.
 \end{aligned}
\]
Here $\ddc_xG=\delta_z-\omega_1$ and $\int G(x,z)\omega_1(x)=0$.
The convolution is continuous and has mean zero.

\begin{theorem}\label{thm:fixed}
For every $\sigma>0$ there is a greatest genuine finite-mass exponential
solution among all $\theta$-psh exponential subsolutions $v\preceq\phi_{a,d}$:
\[
 \MA_\theta(s_\sigma)=e^{\sigma s_\sigma}f\,dW.
\]
It has ordinary bounded type $s_\sigma\sim\phi_{a,d}$ and is unique in that type.
At poles the exponential is zero; no divisor mass is multiplied by an
infinite weight.

Put $v_\sigma=s_\sigma-\sup_Xs_\sigma$ and
$c_\sigma=e^{\sigma\sup_Xs_\sigma}$. As $\sigma\downarrow0$, the entire family
converges in $L^1(X)$ and locally uniformly on $X\setminus D$:
\[
 v_\sigma\longrightarrow v_0=\psi+w_0-S_0,\qquad
 \MA_\theta(v_0)=c_0 f\,dW,\qquad c_\sigma\longrightarrow c_0.
\]
The limit has ordinary type $\phi_{a,d}$, finite positive mass and belongs
to the exact relative stretching class $\E_{\rm env}(X,\theta,\phi_{a,d})$.
The scalar-centered family converges in the same topologies to
\[
 s_0=\psi+w_0-\int_X(\psi+w_0)\,d\nu_0,
 \qquad s_\sigma-\frac{\log c_0}{\sigma}\longrightarrow s_0,
\]
and the actual scalar error satisfies
\[
 \frac1\sigma\log\frac{c_\sigma}{c_0}
 \longrightarrow S_0-\int_X(\psi+w_0)\,d\nu_0.
\]
Within the separated sector $u=\psi+r(x)\preceq\phi_{a,d}$, $s_0$ is
the greatest critical subsolution satisfying
$\MA_\theta(u)\geq c_0f\,dW$ and $\int_Xu\,d\nu_0\leq0$.
The normalized limit is the unique supremum-zero critical member of this sector.
Moreover, $c_0$ is the largest coefficient $c>0$ for which a genuine
$u\in\PSH(X,\theta)$ with $u\preceq\phi_{a,d}$ can satisfy
$\MA_\theta(u)\geq c f\,dW$.
\end{theorem}

\section{Removing the fixed divisor and proving greatestness}

\begin{lemma}\label{lem:removal}
Every genuine $v\in\PSH(X,\theta)$ with $v\preceq\psi$ has a unique representation $v=\psi+r$ with
$r\in\PSH(X,\eta)$. Moreover
\[
 \MA_\theta(\psi+r)=\MA_\eta(r),
\]
and ordinary bounded type relative to $\psi$ is equivalent to bounded $r$.
\end{lemma}
\begin{proof}
Off $D$, the function $r=v-\psi$ is $\eta$-psh and locally bounded above
near $D$. On a coordinate ball choose a strict smooth quadratic $\chi$
with $\ddc\chi\geq\eta$. The plurisubharmonic function $r+\chi$ extends
canonically across $D$ by the bounded-above removable theorem
\cite[I, Theorem 5.24]{Demailly}. Its Hessian is a positive matrix-valued
measure. The measure $\eta+\ddc r$ is positive off $D$; on $D$ it equals
the restriction of $\ddc(r+\chi)$, since smooth corrections assign $D$
zero measure. It is therefore positive everywhere. Local extensions glue
by canonical equality off $D$. On $D$ both $v$ and $\psi+r$ equal $-\infty$,
so their equality holds genuinely there too.

Conversely adding $\psi$ to an $\eta$-psh residual adds the positive current
$d[D]$. Off $D$ the background currents agree. Local bounded expansion with
a common strictly psh correction, followed by genuine finite-level locality,
identifies their non-pluripolar products there. Both products give $D$ zero
mass, so the equality is global, including any additional poles of $r$.
The bounded-type assertion follows from the two comparisons off $D$ and
the canonical extension.
\end{proof}

Set $g_\sigma=e^{\sigma\psi}f$. This is nonzero, nonnegative and continuous,
positive off $D$, with value zero on $D$. Since $\eta$ is smooth positive,
a constant strict barrier is available on all $X$. The bounded exponential
existence theorem \cite[Theorem 4.5]{BGL} gives a bounded solution $R_\sigma$
of
\[
 \MA_\eta(R_\sigma)=e^{\sigma R_\sigma}g_\sigma\,dW.
\]
No bounded mass property is needed in this application of Section 4.
Lemma \ref{lem:removal} makes $\psi+R_\sigma$ an ordinary-type solution.

To prove its greatest property, let $\psi+r\preceq\psi$ be any exponential
subsolution and fix $M>0$. The function $t_M=\max(r,R_\sigma-M)$ is bounded:
the canonical residual has a finite compact upper bound and the floor is
bounded. The local bounded maximum inequality \cite[Lemma 1.2]{GL} gives
\[
 \MA_\eta(t_M)\geq e^{\sigma t_M}g_\sigma\,dW.
\]
For unbounded $r$, apply it to $\max(r,-j)$ with $-j$ below the floor.
On the selected residual branch the genuine finite-level product agrees
with $\MA_\eta(r)$. On the floor branch use
$\MA_\eta(R_\sigma-M)=e^{\sigma R_\sigma}g_\sigma\,dW$.
The ordered maximum inequality includes the equality set.
For $\varepsilon>0$, on $\{R_\sigma<t_M-\varepsilon\}$ we have
\[
 \MA_\eta(R_\sigma)
 \leq e^{-\sigma\varepsilon}\MA_\eta(t_M-\varepsilon).
\]
Both potentials are bounded, their difference is bounded, and a sufficiently
negative constant strict barrier lies below both. The strict domination
principle \cite[Lemma 4.2]{BGL} gives $R_\sigma\geq t_M-\varepsilon$.
Letting $\varepsilon\downarrow0$ proves $R_\sigma\geq r$.
Bounded uniqueness gives uniqueness in ordinary type. This is an order
argument on the full singular class, not merely a separated existence result.

\section{The ordinary model}

First form the envelope on the curve with background $K\omega_1$.
It is genuine and invariant under circle rotations in $z$: the obstacle
is invariant, the family is locally bounded above, and the tangent profile
below supplies a lower competitor for $C\geq C_*$. A common smooth local
correction and the local upper-envelope theorem retain its curvature bound.

With $s=\log|z|^2$, its profile $u+K\log(1+e^s)$ is convex and lies below
$F(s)=K\log(1+e^s)$. The ordinary upper comparison with $\psi$ bounds its
final slope by $a$. The tangent point $F'(s)=a$ is $s=\log(a/d)$.
For $s$ above this point, convexity and the slope bound give
$u+F\leq F(\log R)+a(s-\log R)$; below it the obstacle gives $u\leq0$.
The tangent profile attains both bounds, giving \eqref{eq:model}.
It satisfies $\psi\leq\phi\leq\psi+C_*$ and is a competitor for every
$C\geq C_*$. Thus $P_{K\omega_1}[\psi](0)=\phi$; ordinary-type invariance of
this limit gives the curve model property.
For any full-product competitor $v\leq\min(\psi+C,0)$, its restriction
to $\Pone\times\{y\}$ is either identically $-\infty$ or a genuine
$K\omega_1$-psh curve competitor. Thus $v\leq\phi$ on every slice.
The curve model lifts to a product competitor, since the vertical term is
$b\kappa>0$. Consequently the full product envelope is exactly $\phi$
for every $C\geq C_*$, giving $P_\theta[\phi](0)=\phi$ and $\sup\phi=0$.
No symmetry or K\"ahler assumption on $Y$ is required.

For $\delta>d/(a+d)$, the explicit current
\[
 \delta\theta+\ddc\psi
 =(\delta(a+d)-d)\omega_1+\delta b\kappa+d[D]
\]
is positive. Hence $P_{\delta\theta}(\phi)$ is genuine.
The horizontal non-pluripolar factor of the model is
$K\mathbf1_{\{|z|^2<R\}}\omega_1$, of mass $a$; the tangent derivative
is continuous, so there is no boundary-circle atom. Its actual top measure
is $N b^kK\mathbf1_{\{|z|^2<R\}}dW$, finite and positive.

\section{A contraction giving uniform bounds}

Write $\ddc w=(\Delta w)\omega_1$ and
$C_G=\sup_x\int|G(x,z)|\omega_1(z)<\infty$.
On the continuous mean-zero functions define
\[
 T_\sigma(w)=G*\left(\frac{aH e^{\sigma(\psi+w)}}{Z(w)}-a\right),
 \qquad Z(w)=\int H e^{\sigma(\psi+w)}\omega_1.
\]
Logarithmic convolution of bounded densities is continuous: truncate the
kernel and then use its uniformly integrable logarithmic tail.

The estimates are uniform when $a,d$ range in compact positive intervals
and $h_{\min}\leq H\leq h_{\max}$. Since
$e^{\sigma\psi}\geq1+\sigma\psi$ and $\log\ell$ is integrable,
$\int H e^{\sigma\psi}\omega_1\geq h_{\min}/2$ for small $\sigma$.
On a ball $\|w\|_\infty\leq R_0$, with $\sigma R_0$ small, the normalized
density is bounded by $B=4a_{\max}h_{\max}/h_{\min}$.
Choose $R_0>C_G(B+a_{\max})$.
Then $T_\sigma$ preserves this mean-zero ball. Its density derivative in a
direction $h$ is
\[
 \sigma A(w)\left(h-\frac1a\int h A(w)\omega_1\right),
 \qquad A(w)=\frac{aH e^{\sigma(\psi+w)}}{Z(w)}.
\]
It follows that
\[
 \|T_\sigma(w)-T_\sigma(v)\|_\infty
 \leq2\sigma B C_G\|w-v\|_\infty.
\]
For small $\sigma$ this is a strict contraction. Its mean-zero fixed point
$w_\sigma$ has the same bound $R_0$ for all parameters under consideration.
Set $Z_\sigma=Z(w_\sigma)$ and $A_\sigma=A(w_\sigma)$.
Then
\[
 a+\Delta w_\sigma=A_\sigma,\qquad
 s_\sigma=\psi+w_\sigma+\frac{\log(a/Z_\sigma)}\sigma.
\]
The continuous residual form is
$A_\sigma\omega_1+b\kappa$, positive off $D$ and semipositive on $D$.
Local bounded Bedford--Taylor recursion \cite{BT}, with a common strict
quadratic correction, identifies its actual top product by mollification of the
continuous complex Hessian. Removing $D$ therefore gives
\[
 \MA_\theta(s_\sigma)=N b^k A_\sigma\,dW
                       =e^{\sigma s_\sigma}f\,dW.
\]
Its mass is at most $N a\max b^k$ and is positive.
The preceding greatestness argument identifies it with the full greatest
family. No square of a divisor distribution is taken.

\section{The limits and the order characterization}

The uniform bound and dominated convergence give
$Z_\sigma\to m$ and $A_\sigma\to A_0=aH/m$ in $L^1$.
If $F_j\to0$ in $L^1$ with a uniform bound $B'$, then
\[
 \|G*F_j\|_\infty
 \leq B'\sup_x\int_{\{|G(x,z)|>M\}}|G(x,z)|\omega_1(z)
       +M\|F_j\|_{L^1}.
\]
The first term tends uniformly to zero with $M$. Hence
$w_\sigma\to w_0$ uniformly on the whole curve. In particular
$S_\sigma=\sup(\psi+w_\sigma)\to S_0$ and
$v_\sigma=\psi+w_\sigma-S_\sigma$ has the asserted limits.
Its limiting product is the actual continuous top product
$N b^kA_0dW=c_0 f\,dW$, of finite positive mass.
Since $v_0\geq\phi-C$, for every $t\geq1$ the genuine competitor
$a_t=\phi-tC$ satisfies
$t^{-1}a_t+(1-t^{-1})\phi\leq v_0$.
This is exactly the stretching definition, including common poles.

The actual normalized scalar is $c_\sigma=a e^{\sigma S_\sigma}/Z_\sigma$.
Because $\psi\leq0$ is integrable and $w_\sigma$ is uniformly bounded,
dominated convergence in the exact difference quotient gives
\[
 \frac{Z_\sigma-m}{\sigma}
 =\int H\frac{e^{\sigma(\psi+w_\sigma)}-1}{\sigma}\omega_1
 \longrightarrow\int H(\psi+w_0)\omega_1.
\]
The negative part is bounded by $H(|\psi|+R_0)$ and the positive part
by a fixed constant. Expanding this positive scalar logarithm gives the
centering and scalar-error conclusions of Theorem \ref{thm:fixed}.

For the order characterization take $u=\psi+r(x)\preceq\phi$ with
$\MA_\theta(u)\geq c_0 f\,dW$.
Divisorial removal and the product formula give
$\MA_{a\omega_1}(r)\geq A_0\omega_1$.
Closed curve truncations bound its mass by $a$, which equals the target mass.
Thus the actual measure is exactly $A_0\omega_1$ and has full mass.
The positive closed current $a\omega_1+\ddc r$ itself has mass $a$,
so no measure remains on its pole set. Consequently
$\ddc r=(A_0-a)\omega_1$ as distributions. Harmonic distributions on
$\Pone$ are constant, and canonical equality gives $r=w_0+C$.
The condition $\int u\,d\nu_0\leq0$ is precisely
$C\leq-\int(\psi+w_0)d\nu_0$. This proves the greatest assertion.
Supremum normalization singles out $v_0$. No uniqueness outside this
separated critical sector is needed.

To prove maximality of the coefficient in the full prescribed class, let
$u=\psi+r\preceq\phi$ satisfy $\MA_\theta(u)\geq c f\,dW$, where $c>0$.
Lemma \ref{lem:removal} supplies the genuine residual $r$ and its finite
upper bound $B=\sup_Xr$. For every $\sigma>0$ the potential
\[
 u+\frac{\log c}{\sigma}-B
\]
is an exponential subsolution in this class, since
$e^{\sigma(\psi+r-B)}\leq1$. Greatestness therefore gives
$r+\sigma^{-1}\log c-B\leq q_\sigma$ off $D$, where
$s_\sigma=\psi+q_\sigma$. Fix a point off $D$ where $r$ is finite.
The already proved centering limit bounds
$q_\sigma-\sigma^{-1}\log c_0$ there. Multiplying the inequality by
$\sigma$ and letting $\sigma\downarrow0$ gives $\log c\leq\log c_0$.
The potential $v_0$ attains $c_0$. This argument permits additional poles
and infinite competitor mass; it asserts no uniqueness of nonseparated
critical potentials.

\section{Moving ordinary references at independent rates}

\begin{theorem}\label{thm:moving}
Let $a_j,d_j>0$, $a_j\to a>0$, $d_j\to d>0$, $b_j\to b$ uniformly, with each $b_j$
smooth positive and $b$ smooth positive nonconstant. Let continuous positive
$H_j\to H$ uniformly. For any $\sigma_j>0$ tending to zero, set
\[
 \begin{aligned}
 \theta_j&=(a_j+d_j)\omega_1+b_j\kappa,
 &\psi_j&=d_j\log\ell,\\
 \phi_j&=\phi_{a_j,d_j},&f_j&=N b_j^kH_j,\\
 m_j&=\int H_j\omega_1,&c_{0,j}&=a_j/m_j.
 \end{aligned}
\]
The normalized greatest solutions tend in $L^1(X)$ and locally uniformly
off $D$ to $v_0$ from Theorem \ref{thm:fixed}. Their actual measures tend in
total variation to $c_0 f\,dW$. The limit has ordinary bounded type
$\phi_{a,d}$ and the exact relative stretching class.
Writing $c_j=e^{\sigma_j\sup s_j}$, one has
\[
 c_j\to c_0,\qquad
 \frac1{\sigma_j}\log\frac{c_j}{c_{0,j}}
 \longrightarrow S_0-\int(\psi+w_0)d\nu_0,\qquad
 s_j-\frac{\log c_{0,j}}{\sigma_j}\longrightarrow s_0.
\]
There is no restriction on the rates of parameter and exponent variation.
\end{theorem}
\begin{proof}
The contraction bounds hold uniformly for all late parameters. Their
mean-zero fixed points $w_j$ are uniformly bounded. Dominated convergence
gives $Z_j\to m$ and $A_j\to A_0$ in $L^1$. The bounded-density Green
estimate then gives $w_j\to w_0$ uniformly. At the pole $A_j$ equals zero,
whereas $A_0$ is strictly positive; no uniform density convergence is claimed.

The bounds for $w_j$ and $d_j\geq d_{\min}>0$ keep all maximizers of
$\psi_j+w_j$ uniformly away from $D$: a point with $\ell=1$ supplies a
uniform competing value, while $d_{\min}\log\ell+R_0$ tends uniformly to
$-\infty$ near $D$. On the compact complement $\psi_j\to\psi$ uniformly,
so $S_j\to S_0$. This proves the normalized limits, including $L^1$ because
$\log\ell$ is integrable. Uniform convergence of $b_j$ and $L^1$ convergence
of $A_j$ give actual total-variation convergence of $N b_j^k A_jdW$.

In \eqref{eq:model}, $R_j\to R$ and $C_{*,j}\to C_*$. The bounded continuous
residual extensions of $\phi_j-\psi_j$ across $D$ converge uniformly, including the moving tangent
circle and infinity: on the two pieces they equal respectively
$d_j\log(1+|z|^2)$ and
$a_j\log(|z|^2/(1+|z|^2))+C_{*,j}$.
The tangent radii stay in a compact positive interval and the pieces agree
there. In particular the uniform bound $0\leq\phi_j-\psi_j\leq C_{*,j}$
and the explicit limit prove $v_0\sim\phi_{a,d}$. This establishes actual
limiting type, rather than presuming continuity of singularity types.

Finally use the difference quotient
\[
 \frac{Z_j-m_j}{\sigma_j}
 =\int H_j\frac{e^{\sigma_j(\psi_j+w_j)}-1}{\sigma_j}\omega_1
 \longrightarrow\int H(\psi+w_0)\omega_1.
\]
It has the common integrable negative bound supplied by
$d_{\max}|\log\ell|+R_0$ and a constant positive bound. Since
$c_j/c_{0,j}=m_j e^{\sigma_j S_j}/Z_j$, the asserted scalar error and
centering follow. Replacing $c_{0,j}$ by $c_0$ would add
$\sigma_j^{-1}\log(c_{0,j}/c_0)$, which need not vanish.
\end{proof}

For a concrete nonclosed example take $a=1$, $d=1/4$, $b=2+t$ and $H=1$,
where $t$ is the standard sphere height. Then $w_0=0$, $S_0=0$,
$\int\log\ell\,\omega_1=-1$, $c_0=1$, and the selected centered limit is
$\psi+1/4$, with scalar-error limit $1/4$.
For moving parameters $a_j=1+\sqrt{\sigma_j}$ and fixed $d,b,H$,
the same theorem centers by $c_{0,j}=a_j$. Centering by the limiting scalar
$1$ adds $\log(1+\sqrt{\sigma_j})/\sigma_j\to+\infty$.
The pole remains present in every solution and in the limiting ordinary type.
For an actual non-K\"ahler product one may instead take $Y$ to be the scalar
Hopf surface $(\mathbf C^2\setminus\{0\})/\langle z\mapsto2z\rangle$ with
the descended positive form $i\sum dz_j\wedge d\bar z_j/|z|^2$, rescaled
to have unit square integral. It is diffeomorphic to $S^1\times S^3$ and has
first Betti number one, so it is not K\"ahler. A K\"ahler product would restrict
to a K\"ahler metric on this compact fiber, which is impossible.

\begin{thebibliography}{9}
\bibitem{BT} E. Bedford and B. A. Taylor,
\emph{A new capacity for plurisubharmonic functions}, Acta Math. 149 (1982), 1--40.
\bibitem{BGL} S. Boucksom, V. Guedj and C. H. Lu,
\emph{Volumes of Bott--Chern classes}, arXiv:2406.01090v1.
\bibitem{GL} V. Guedj and C. H. Lu,
\emph{Quasi-plurisubharmonic envelopes 2: Bounds on Monge--Amp\`ere volumes},
arXiv:2106.04272.
\bibitem{Demailly} J.-P. Demailly,
\emph{Complex analytic and differential geometry}, freely available manuscript,
I, Section 5.F.
\end{thebibliography}
\end{document}
